\documentclass[11pt]{article}
\usepackage[margin=1in]{geometry}
\usepackage{amsmath,amssymb,booktabs,graphicx,microtype}
\usepackage[colorlinks=true,linkcolor=blue!50!black,citecolor=blue!50!black]{hyperref}
\usepackage{caption}
\captionsetup{font=small}
\title{\textbf{Can the inventory kernel be identified?}\\[2pt]
\large A simulation study: data-generating process, estimator, and results}
\author{ProductGPT technical note (R34)}
\date{28 September 2026}

\begin{document}
\maketitle

\begin{abstract}
\noindent The model we are building carries a \emph{stock} of previously obtained
products that decays over time and spreads across similar products. Three
objects have to be estimated: how strongly one product's acquisition satiates
demand for another ($\kappa$), how fast that effect fades ($\varphi$), and how
much a duplicate copy counts ($g$). This note asks which of them our data can
identify, using a simulation in which the truth is known by construction.
The answer: $\varphi$ and $g$ are recoverable, $\kappa$ is not, and the reason
is the rotating assortment rather than the estimator. A parametric kernel with
two coefficients partially survives; a free kernel does not and, worse, invents
structure when none exists. Section~\ref{sec:real} reports what happened when
the resulting model met the real data.
\end{abstract}

\section{The question, and why a simulation answers it}

On real data we can only ever compare fits. If a flexible kernel improves the
likelihood we cannot tell whether it found substitution or absorbed something
else; if it does not improve the likelihood we cannot tell whether substitution
is absent or merely invisible. A simulation removes the ambiguity: we generate
data from a known kernel, fit the estimator, and compare the estimate with the
truth. Everything below is synthetic. No consumer data is read.

\section{The data-generating process}\label{sec:dgp}

\subsection{Objects}

$P$ products, each carrying an attribute (its ``element'') $e(j) \in
\{1,\dots,E\}$. $N$ customers, $T$ decision occasions each. At every occasion
$B$ banners are live simultaneously; banner $b$ features $F$ products, so $BF$
distinct products are on offer at once. The assortment is fixed within a
campaign of $L$ occasions and rotates between campaigns.

\subsection{The state}

Let $w_\tau(p)$ be the weight of an acquisition of product $p$ at occasion
$\tau$. The customer's stock of product $j$ entering occasion $t$ is
\begin{equation}
S_t(j) \;=\; \sum_{\tau < t} \kappa(j, p_\tau)\; g(k_\tau)\; \varphi(t-\tau),
\qquad \varphi(\Delta) = \rho^{\Delta},\quad \rho = 2^{-1/h},
\label{eq:stock}
\end{equation}
where $h$ is the half-life in occasions, $k_\tau$ is which copy of $p_\tau$ this
was ($1$st, $2$nd, $3$rd), and $g$ is weakly decreasing with $g(1)=1$. The
kernel is row-stochastic,
\begin{equation}
\kappa(j,p) \;\propto\; \exp\big(\theta_{\text{self}}\,\mathbf{1}[j=p]
   \;+\; \theta_{\text{attr}}\,\mathbf{1}[e(j)=e(p)]\big),
\label{eq:kappa}
\end{equation}
so $\theta_{\text{attr}}=0$ means no substitution across products and
$\kappa = I$.

\subsection{The decision}

Nine alternatives: $B$ banners $\times$ $\{$one draw, ten draws$\}$, plus
\emph{NotBuy}. Banner $b$'s attractiveness averages taste net of satiation over
the products it offers,
\begin{equation}
A_{ib t} \;=\; \frac{1}{F}\sum_{j \in \mathcal{A}_{bt}}
   \big(\theta_i(j) - \lambda_i\, S_{it}(j)\big),
\end{equation}
with $\theta_i(j) = \text{base}(j) + \text{taste}_i(e(j))$ and
$\text{taste}_i \sim \mathcal{N}(0,\sigma_{\text{taste}}^2)$. Utilities are
\begin{equation}
V_{i,(b,1),t} = c_1 + \beta A_{ibt} + \xi_{c(t)}, \qquad
V_{i,(b,10),t} = c_{10} + \beta A_{ibt} + \xi_{c(t)}, \qquad
V_{i,\text{NotBuy},t} = 0,
\end{equation}
where $\xi_{c}$ is a campaign-level demand shock with standard deviation
$\sigma_{\xi}$. The choice is multinomial logit (Gumbel errors).

\subsection{Resolution of a purchase}

If the customer buys $n \in \{1,10\}$ draws on banner $b$, each draw returns the
banner's featured 5-star with probability $p_5$, its featured 4-star with
probability $p_4$, and otherwise a uniformly drawn filler. Composition is
therefore \emph{exogenous} given the number of draws --- which product you
receive is the machine's decision, not yours. This matters in
Section~\ref{sec:theory}.

\subsection{Parameters}

\begin{table}[h]\centering\small
\begin{tabular}{llll}
\toprule
Symbol & Meaning & Baseline & Varied in \\
\midrule
$P,E$ & products, attributes & 24, 6 & --- \\
$N,T$ & customers, occasions & 1000, 300 & s1 ($N$) \\
$B,F$ & concurrent banners, featured each & 4, 2 & --- \\
$L$ & occasions per campaign & 15 (45 in s6--s7) & s4 \\
$h$ & half-life (occasions) & 30 & --- \\
$\theta_{\text{self}},\theta_{\text{attr}}$ & kernel coefficients & 3.0, 1.5 & s2, s6 \\
$g$ & duplicate weights & $(1,\,0.6,\,0.35)$ & --- \\
$\lambda$ & satiation strength & 1.0 & s0, s7 \\
$\sigma_{\text{taste}}$ & taste heterogeneity & 0 & s2 \\
$\sigma_{\lambda}$ & heterogeneity in satiation & 0 & s7 \\
$\sigma_{\xi}$ & campaign demand shock & 0 & s6 \\
$p_5,p_4$ & gacha rates & 0.01, 0.10 & --- \\
\bottomrule
\end{tabular}
\end{table}

\section{The estimator}

The estimator has the same functional form as \eqref{eq:stock}--\eqref{eq:kappa}
and is fitted by maximum likelihood (Adam, minibatches over customers, 600
epochs). Free parameters: the kernel, $\log h$, the duplicate steps, $\lambda$,
$\beta$, the two intercepts, and a per-product base taste. Three kernel
specifications are compared:

\begin{center}\small
\begin{tabular}{lll}
\toprule
Specification & Free parameters & Interpretation \\
\midrule
\texttt{identity} & 0 & no substitution across products \\
\texttt{attr} & 2 & McAlister attribute satiation, coefficients estimated \\
\texttt{learned} & $P^2 = 576$ & a free QKV attention kernel \\
\bottomrule
\end{tabular}
\end{center}

\noindent Two options are switched on where indicated: a \emph{customer fixed
effect} (per-customer attribute tastes, the standard remedy for heterogeneity)
and a \emph{campaign fixed effect} (a per-campaign intercept on the purchase
alternatives).

The first 20 occasions are burn-in. Where a held-out figure is reported, the
model is fitted on the first 70\% of occasions and scored on the rest.

\paragraph{Metrics.} Recovery of $\kappa$ is measured by the Spearman rank
correlation between the off-diagonal entries of $\hat\kappa$ and $\kappa^{*}$,
which is invariant to the scale and to any monotone distortion. ``Element
lift'' is the ratio of mean off-diagonal mass within an attribute group to that
outside it; $1.0$ means no structure was found. For the parametric kernel we
compare $\hat\theta$ with $\theta^{*}$ directly.

\paragraph{One implementation detail that mattered.} Copies acquired in the
same ten-draw pull must be weighted as the 1st, 2nd and 3rd copy. An earlier
version of the estimator gave them all the pre-occasion holding, which
over-weighted duplicates; it was caught by replaying the state against an
explicit evaluation of \eqref{eq:stock} and is fixed.

\section{Results}

\subsection{s0 --- how much is the kernel worth at all?}

Before asking whether $\kappa$ can be estimated, ask whether it matters. We fit
two models at equal budget: one with $\kappa$ frozen at the \emph{truth}, one
with $\kappa = I$. The gap is an upper bound on what any estimator could gain.

\begin{center}\small
\begin{tabular}{lrrr}
\toprule
$\lambda$ & identity kernel & true kernel & gap \\
\midrule
0.5 & 1.5910 & 1.5910 & $0.000$ \\
1.0 & 1.4330 & 1.4320 & $0.000$ \\
2.0 & 1.1960 & 1.1960 & $0.000$ \\
4.0 & 0.9020 & 0.9010 & $+0.001$ \\
\bottomrule
\end{tabular}
\end{center}

\noindent Knowing the true kernel is worth at most $0.001$ nats, at every
satiation strength. Note that satiation \emph{itself} matters enormously over
this range --- the NotBuy share moves from 0.49 to 0.76 --- so the flat quantity
is specifically the \emph{cross-product structure}, not the stock.

\subsection{s1 --- recovery, and how many customers it would take}

\begin{center}\small
\begin{tabular}{lrrrr}
\toprule
$N$ & Spearman($\hat\kappa,\kappa^{*}$) & element lift & $\hat h$ (true 30) & $\hat\lambda$ (true 1.0) \\
\midrule
250 & 0.062 & 1.90 & 30.7 & 0.88 \\
500 & 0.013 & 1.94 & 28.5 & 0.61 \\
1000 & 0.080 & 1.64 & 29.7 & 0.89 \\
2000 & 0.099 & 0.38 & 30.8 & 0.75 \\
\bottomrule
\end{tabular}
\end{center}

\noindent The half-life is recovered at every sample size. The kernel is not,
and \emph{more customers do not help} --- which already indicates a structural
problem rather than a power problem.

\subsection{s2 --- does the kernel invent structure?}

Here the truth is $\kappa^{*} = I$: no substitution whatsoever. A well-behaved
estimator should return a near-diagonal matrix.

\begin{center}\small
\begin{tabular}{llrrr}
\toprule
DGP & customer FE & element lift & off-diagonal mass & diagonal \\
\midrule
homogeneous & no & 1.48 & 0.46 & 0.54 \\
homogeneous & yes & 1.54 & 0.66 & 0.34 \\
heterogeneous tastes & no & 1.08 & 0.60 & 0.40 \\
heterogeneous tastes & yes & \textbf{2.12} & 0.70 & 0.31 \\
\bottomrule
\end{tabular}
\end{center}

\noindent The free kernel places 46--70\% of its mass off the diagonal when the
truth is the identity, and the customer fixed effect --- the textbook remedy ---
makes it \emph{worse}. A flexible estimator does not report ``I don't know''; it
reports a confident matrix.

\subsection{s4 and s6(A) --- the offer schedule is the binding constraint}

\begin{center}\small
\begin{tabular}{lrr}
\toprule
Offers rotate\dots & Spearman & $\hat h$ \\
\midrule
every 30 occasions & $-0.015$ & 29.2 \\
every 15 occasions & $0.007$ & 29.6 \\
every 5 occasions & $0.072$ & 28.8 \\
randomised every occasion & $\mathbf{0.648}$ & 29.0 \\
\bottomrule
\end{tabular}
\end{center}

\noindent The estimator recovers $\kappa$ when offers are randomised and fails
at every realistic rotation speed. This is a \emph{cliff, not a slope}:
rotating six times faster than the baseline buys almost nothing. So the
estimator is sound and the design is the constraint.

Fixing the schedule at our own calendar ($L = 45$) and varying the true
substitution strength (s6A) shows how strong the effect would have to be:

\begin{center}\small
\begin{tabular}{lrr}
\toprule
$\theta_{\text{attr}}^{*}$ & $\hat\theta_{\text{attr}}$ & \% of truth \\
\midrule
0.5 & 0.76 & 152\% \\
1.5 & 0.82 & 55\% \\
3.0 & 2.42 & 81\% \\
5.0 & 2.81 & 56\% \\
\bottomrule
\end{tabular}
\end{center}

\noindent The estimate has a \emph{floor} near $0.8$ whatever the truth: below
$\theta^{*}\approx 3$ it reports roughly the same number regardless, so a
weak-but-real effect and a biased constant are indistinguishable.

\subsection{s5 --- structure buys identification, partially}

Replacing the free kernel by the two-coefficient parametric form
\eqref{eq:kappa}:

\begin{center}\small
\begin{tabular}{llrrrr}
\toprule
Schedule & kernel & free params & held-out NLL & $\hat\theta_{\text{self}}$ (3.0) & $\hat\theta_{\text{attr}}$ (1.5) \\
\midrule
campaign rotation & identity & 0 & 1.6062 & --- & --- \\
campaign rotation & parametric & 2 & 1.6060 & 3.28 & 0.89 \\
campaign rotation & free & 576 & 1.6063 & --- & --- \\
randomised & identity & 0 & 1.5718 & --- & --- \\
randomised & parametric & 2 & 1.5715 & 3.05 & 1.40 \\
randomised & free & 576 & 1.5776 & --- & --- \\
\bottomrule
\end{tabular}
\end{center}

\noindent Under randomisation the parametric kernel is recovered essentially
perfectly (rank correlation 1.00). Under a rotation its own-product coefficient
survives but the cross-attribute coefficient is attenuated by about 40\%.
Note also that the 576-parameter kernel is \emph{not} better out of sample than
assuming no substitution at all.

\subsection{s6(B) --- campaign shocks masquerade as memory}

Time since acquisition and time since the product left the assortment advance
together. Without a campaign control the decay absorbs the calendar:

\begin{center}\small
\begin{tabular}{lrr}
\toprule
$\sigma_\xi$ & $\hat h$, no campaign FE & $\hat h$, with campaign FE \\
\midrule
0.0 & 30.4 & 30.2 \\
0.3 & \textbf{59.3} & 29.5 \\
0.6 & \textbf{147.2} & 30.4 \\
\bottomrule
\end{tabular}
\end{center}

\noindent True $h = 30$. A modest calendar shock doubles the estimate and a
larger one quintuples it. \textbf{The empirical model must carry campaign fixed
effects or its half-life is not interpretable.}

\subsection{s7 --- heterogeneous satiation against a homogeneous estimator}

Each customer is given their own $\lambda_i$ (log-normal); the estimator fits a
single $\lambda$.

\begin{center}\small
\begin{tabular}{lrrr}
\toprule
$\sigma_\lambda$ & population mean $\lambda$ & $\hat\lambda$ & shortfall \\
\midrule
0.0 & 1.00 & 0.93 & $-7\%$ \\
0.3 & 1.04 & 0.77 & $-26\%$ \\
0.6 & 1.19 & 0.77 & $-36\%$ \\
1.0 & 1.61 & 0.74 & $-54\%$ \\
\bottomrule
\end{tabular}
\end{center}

\noindent A pooled satiation estimate is a \emph{lower bound}. Unlike s2's
preference heterogeneity, slope heterogeneity does not manufacture false
substitution --- but a customer fixed effect again inflates the kernel rather
than curing it.

\section{Why rotation blocks identification}\label{sec:theory}

Under a block rotation, acquisitions can only come from the current assortment
$\mathcal{A}_c$, so \eqref{eq:stock} aggregates:
\begin{equation}
S_{it}(j) \;\approx\; \sum_{c' \le c} \varphi(t - T_{c'})\; N_{i,c'}\;
   K(j, \mathcal{A}_{c'}),
\qquad
K(j,\mathcal{A}) = \sum_{p \in \mathcal{A}} \kappa(j,p)\, q_p ,
\label{eq:collapse}
\end{equation}
where $N_{i,c'}$ is how much customer $i$ pulled in campaign $c'$ and $q_p$ is
$p$'s gacha rate. Two institutional facts produce this: composition within a
banner is exogenous, and everything in an assortment arrives in the same window
and so shares one decay factor.

Three consequences follow.

\begin{enumerate}
\item \textbf{Only assortment-level aggregates enter the likelihood.} The
contrast $\kappa(j,p_1) - \kappa(j,p_2)$ for two products offered together is
identified only by gacha sampling noise. With 49 products over 30 campaigns
there are $49 \times 48 = 2352$ off-diagonal unknowns against at most
$49 \times 30 = 1470$ weakly-informed functionals.

\item \textbf{Cross terms are attenuated before they are seen.} Our campaigns
run a median of 45 occasions, so a cross-product effect is first observable
after a delay of that order: a factor $2^{-45/h}$, about $0.35$ at $h=30$.
Own-product terms suffer no such penalty because $j$ is offered throughout its
own campaign.

\item \textbf{The surviving variation is collinear with the propensity.}
$N_{i,c'}$ multiplies every term, so cross-product stock is nearly proportional
to how much the customer recently pulled --- which own-stock, previous decision
and any customer effect already absorb. This is why the customer fixed effect
made s2 worse rather than better.
\end{enumerate}

\noindent $\varphi$ and $g$ escape all three because they rest on different
contrasts: $\varphi$ on \emph{within-customer timing} (the same customer pulls
early and late within a campaign), and $g$ on \emph{within-product} variation
(holding one copy versus three), which occurs inside a single assortment and is
driven by exogenous gacha luck.

\section{What happened on the real data}\label{sec:real}

The model implied by these results --- per-product counts, learned decay,
duplicate weights, an attribute kernel, campaign fixed effects --- was fitted to
the Genshin panel (5 seeds per arm).

\paragraph{Duplicate weights are estimated.} $\hat g = (1.00,\,0.42 \pm 0.03,\,
0.18 \pm 0.015)$: a second copy is worth about 40\% of the first, a third about
18\%. Stable across seeds and across three different half-life initialisations.

\paragraph{The half-life is not estimable by gradient descent, but the
likelihood is not flat in it.} Starting the optimiser at 10, 30 and 100 returns
9.9, 29.8 and 99.7 --- the parameter never moves. Yet the fit depends on the
value:

\begin{center}\small
\begin{tabular}{lrr}
\toprule
Half-life (fixed at) & holdout NLL & \\
\midrule
10 & $0.8845 \pm 0.0065$ & best \\
30 & $0.8914 \pm 0.0058$ & \\
100 & $0.8946 \pm 0.0098$ & \\
\bottomrule
\end{tabular}
\end{center}

\noindent The gradient is too small to move a parameter of that magnitude while
the likelihood still has a slope along it. The appropriate estimator is
therefore a \textbf{profile likelihood}: freeze $h$, fit everything else, and
read the curve. That grid (half-lives 2, 5, 10, 20, 40) is running; the optimum
appears to lie below 10 occasions.

\paragraph{Caveat on transferring s1's optimism.} The simulation's satiation was
far stronger than the real data appear to carry --- there, removing the entire
stock path costs only 0.0137 nats. So ``the half-life is identified'' held in
simulation at $\lambda=1$ and should not be assumed at the real data's effective
signal strength. The profile likelihood is what settles it.

\section{Implications}

\begin{enumerate}
\item \textbf{Report} the duplicate weights, and the half-life only as a profile
curve with its uncertainty --- never as a point estimate from gradient descent.
\item \textbf{Carry campaign fixed effects} in any specification that reports a
decay.
\item \textbf{Report the pooled satiation as a lower bound} until a hierarchical
version exists.
\item \textbf{Do not report the learned kernel as substitution.} It may appear
as a specification bound: if it cannot beat the parametric restriction out of
sample, nothing beyond the restriction is in the data.
\item \textbf{The negative result is a contribution.} A learned attention
inventory identifies the decay of consumption but not the substitution
structure, under a live-service rotation schedule --- with the randomised-offer
arm proving the estimator is sound and the design is the binding constraint.
It also says what experiment would fix it: randomise one banner slot.
\end{enumerate}

\vspace{1em}
\noindent\footnotesize Code: \texttt{simulation/kernel\_sim.py} (process and
estimator), \texttt{simulation/run\_sim\_study.py} (experiments s0--s7),
\texttt{scripts/test\_kernel\_sim.py} (16 correctness checks, including a
brute-force replay of the state recursion). Full run records and the bugs found
along the way are in \texttt{EXPERIMENTS.md}.

\end{document}
